\section{Economía y estrategias de combinación de modelos}

\begin{importantbox}{Principio de combinación de modelos fuertes y débiles}
\begin{itemize}
    \item \textbf{Planning y RePlan}: se usa un modelo fuerte (como un modelo de razonamiento), porque la planeación determina el límite superior de todo el Workflow
    \item \textbf{Execute}: se usa un modelo débil (como GPT-4.1 Nano/Mini), porque ya se le indicó al modelo qué hacer y solo necesita ejecutarlo paso a paso
\end{itemize}
Esto coincide por completo con el posicionamiento que hace OpenAI de los modelos de razonamiento y los que no lo son.
\end{importantbox}

Si la etapa de Planning es especialmente importante, también se puede introducir el patrón \textbf{Generator-Critique} para formar un Plan más sólido y confiable; no hay que esperar que el modelo proponga un buen Plan en un solo intento.

\subsection{Resumen de la sección}

Una estrategia razonable de combinación de modelos permite equilibrar el costo y el resultado: los modelos caros se encargan de la planeación y los modelos baratos, de la ejecución.

